\scanpage{24}

est nul ; or

\[q^*(e_{\unsure{j|i},m})=m e_i,\qquad p_*(e_{i,n})=\unsure{n e_i|mn e_i}\]

où $e_i$ définit l'extension $\Phi_i$. Donc la \unsure{condition|question} est que

\[(qm-pn)e_i=0\]

On fait l'hypothèse

\[\text{\textcircled{c}}\quad\boxed{e_i\ \text{est une base de }\operatorname{Ext}^{1}_{C_0}(\Psi^i,\Psi^{i-1})}\]

de sorte que la condition \unsure{désirée|ci-dessus|devient} est

\[ qm=pn \]

\[ \text{i.e.}\qquad \frac qp=\frac{\unsure{m|n}}{\unsure{n|m}}=\frac{m'}{n'} \]

si $\unsure{m|m'}=\frac{m}{(m,n)}$, $n'=\frac{n}{(m,n)}$, donc on aura

\[p=\deleted{\unsure{p',\,q}}\,t\unsure{n'|m'},\quad q=tm'\qquad (t\in k)\]

Or pour $p,q$ \unsure{correspondant|correspondants} ainsi à un $t$ (\unsure{évidemt|d'ailleurs} uniquement déterminé par $p,q$) l'hom $\alpha$ qui leur correspond est déterminé mod un $\deleted{\unsure{\operatorname{Ext}|\operatorname{Ex}}}\operatorname{Hom}(\Psi^i,\Psi^{i-1})=0$, donc est unique. On trouve donc

\[ \operatorname{Hom}(\Phi(i,n),\Phi(i,m))\simeq k \]

% archive sheet number: 23
